\section{1. Introduction}\label{ddjf-1}

\subsection{1.1. Purpose of the Document}\label{ddjf-1-1}

The Design Justification File (DDJF) presents the justification of the design solutions adopted for the Ground Segment (GSEG, CI-02) of the Early Warning Alarm System (EWAS).

The document is intended to demonstrate the extent to which the design defined in the GSEG Design Definition File (DDF) satisfies the requirements of the GSEG Technical Requirements Specification (TRS). The justification comprises technical reasoning, option analyses, calculations, margin assessment and available verification evidence.

The document was prepared in accordance with ECSS-E-ST-10C Rev.1 (15 February 2017), Annex K.

\subsection{1.2. Scope of the Document}\label{ddjf-1-2}

The DDJF covers the functional and physical architecture of integrated GSEG, allocation of functions among the seven subsystems, interfaces, technical budgets, service availability and constraints imposed by the design.

Detailed internal solutions for individual subsystems belong to lower-level documentation. SSEG space elements and recipient systems beyond the EWAS interface boundaries are outside the scope of GSEG design justification.

\subsection{1.3. Justification Principles}\label{ddjf-1-3}

Each justification is related to a TRS requirement and the corresponding DDF solution.

Values originating from baseline requirements are distinguished from proposed design allocations. Compliance is considered demonstrated only to the extent supported by analysis, test, review of design or inspection.

The DDJF is updated as the design develops. Lack of appropriate evidence means that the justification remains open.
